\section{案例拆解：一次 Reasoning 升级实验如何设计}

\subsection{实验目标与约束}
假设我们要把一个 7B 级模型在数学与代码推理任务上提升 5--10 个点，同时不明显损伤通用对话能力。预算约束是 ``训练可控、推理可上线''。这类问题正好对应本讲提出的组合路线。

\begin{importantbox}{目标设定方式}
目标应同时包含：
\begin{itemize}
    \item \textbf{能力目标}：可验证任务准确率提升；
    \item \textbf{副作用上限}：通用任务下降不超过阈值；
    \item \textbf{部署目标}：成本和延迟在可接受区间。
\end{itemize}
缺一项都会让优化方向偏离真实需求。
\end{importantbox}

\subsection{分阶段实验矩阵}
\begin{table}[H]
\centering
\renewcommand{\arraystretch}{1.2}
\begin{tabular}{p{2.8cm}p{3.4cm}p{3.8cm}p{3.0cm}}
\toprule
\textbf{阶段} & \textbf{变量} & \textbf{观测指标} & \textbf{停止条件} \\
\midrule
SFT & 人类/合成数据配比、长度分布 & 指令遵循、格式稳定性、基础正确率 & 副作用超过阈值即回滚 \\
DPO/PPO & 偏好损失权重、采样温度 & 有用性、冗长度、拒答率 & 偏好收益边际下降 \\
RLVR & 验证器规则、采样轮数、更新步数 & 数学/代码可验证通过率 & 出现奖励捷径信号 \\
Inference-time & N 值、重排策略、检索轮次 & 通过率、P95 延迟、单位请求成本 & 成本超预算 \\
\bottomrule
\end{tabular}
\caption{Reasoning 升级实验的分阶段矩阵}
\end{table}

\subsection{错误分析优先级}
\begin{knowledgebox}{最有价值的三类错误样本}
\begin{enumerate}
    \item \textbf{高置信错误}：模型给出自信但错误的长链推理；
    \item \textbf{验证器误判}：答案正确但被规则误杀，或相反；
    \item \textbf{预算不经济样本}：需要大量采样才提升极少准确率。
\end{enumerate}
优先修复这些样本，通常能同时提升质量与成本效率。
\end{knowledgebox}

\begin{warningbox}{不要用平均分掩盖结构性失败}
平均分提升可能掩盖关键任务退化。真实系统更需要看 ``任务分桶表现''：例如多步数学、长代码修复、检索依赖问答分别变化如何。
\end{warningbox}

\subsection{把案例映射回本讲主线}
这个案例表明，本讲的三阶段框架和开放配方思想可以直接转化为工程决策：
\begin{itemize}
    \item 预训练决定底座上限；
    \item 后训练决定行为方向；
    \item 推理时扩展决定部署侧收益曲线；
    \item 开放与可复现决定团队能否持续迭代。
\end{itemize}

\subsection{本章小结}
一次成功的 reasoning 升级实验，不是某个单点参数的胜利，而是目标约束、阶段矩阵、错误分析和预算管理协同工作的结果。


